\documentclass[11pt]{article}
\usepackage[margin=1in]{geometry}
\usepackage[T1]{fontenc}
\usepackage[utf8]{inputenc}
\usepackage{lmodern,microtype,amsmath,amssymb,amsthm,mathtools,needspace}
\usepackage[colorlinks=true,linkcolor=blue,citecolor=blue,urlcolor=blue]{hyperref}
\hypersetup{pdftitle={Exponential sample lower bounds for sparse parity under a globally separated neural prior},pdfauthor={Dmitry Vaintrob; Kaarel Hanni; OpenAI Codex}}
\setlength{\emergencystretch}{2em}
\newtheorem{theorem}{Theorem}[section]
\newtheorem{lemma}[theorem]{Lemma}
\newtheorem{proposition}[theorem]{Proposition}
\newtheorem{corollary}[theorem]{Corollary}
\theoremstyle{remark}\newtheorem{remark}[theorem]{Remark}
\DeclareMathOperator{\Var}{Var}
\DeclareMathOperator{\sign}{sign}
\DeclareMathOperator{\dist}{dist}
\DeclareMathOperator{\KL}{KL}
\newcommand{\E}{\mathbb E}
\newcommand{\Pp}{\mathbb P}
\title{Exponential sample lower bounds for sparse parity under a globally separated neural prior}
\author{Dmitry Vaintrob \and Kaarel H\"anni \and OpenAI Codex}
\date{October 15, 2026}
\begin{document}
\maketitle
\begin{abstract}
We study a fixed-depth Bayesian neural network with Gaussian weights,
hyperbolic-tangent hidden activations, and a linear output, conditioned
to have a fixed absolute output margin on every Boolean input and then
on noiseless training labels. For every fixed width-scaling exponent
$0<\gamma<1$, we prove an $\exp(\Omega(k))$ sample lower bound for
learning parity on any $k$ coordinates of a $d$-dimensional input,
uniformly for growing $k\le d$. With probability at least $3/4$ over
the training sample, every sign function selected in the width-first
limit has test error at least $3/8$. Consequently, at least $1/16$ of
inputs lack correct-label posterior probability greater than $2/3$,
regardless of the actual mixture weights among tied minimizers.
A random-restriction argument shows that a low-degree polynomial
approximating a globally separated accurate parity predictor must have
exponentially large squared amplitude in the ratio of parity degree to
polynomial degree. We combine this with the paired Fourier spectrum of
the limiting kernel and uniform bounds on sampled tail covariance
matrices. A second argument improves the sample range to
$\exp\{c\max[k,\sqrt{k}\log(e(d+1)/\sqrt{k})]\}$ in very large ambient
dimensions. All width-limit, spectral, sampling, and interpolation
arguments are included. A generalization upper bound has scale
$\exp\{O(k\log(ed/k))\}$ for growing sparse $k$. This rules out a
literal $d^{\Omega(k)}$ lower bound uniformly over all $k=o(d)$;
a lower bound on the natural $\exp\{\Omega(k\log(ed/k))\}$ scale
remains open.
\end{abstract}

\section{Model and main result}
Let $\mathcal X=\{-1,1\}^d$ with uniform probability measure.
Fix $U\subseteq[d]$, $|U|=k$, and let
$f(x)=\chi_U(x)=\prod_{j\in U}x_j$.
The learner is not given $U$. All hidden layers have width $N$, the
number of weight matrices $L\ge2$ is fixed, and
\[
 h^0(x)=(1,x),\qquad
 h^s(x)=\tanh(W^s h^{s-1}(x))\quad(1\le s<L),\qquad
 z_N(x)=W^L h^{L-1}(x).
\]
There are no further biases. All weights are independent centered
Gaussians with variances
\begin{equation}\label{eq:weights}
 \Var W^1_{ij}=N^{1-\gamma}/d,\qquad
 \Var W^s_{ij}=N^{-\gamma}\quad(1<s<L),\qquad
 \Var W^L_i=N^{-\gamma-1}.
\end{equation}
Fix $a>0$, for example $a=1/2$, and a fixed $0<\gamma<1$.
Condition the prior on $|z_N(x)|>a$ for every input $x$.
For independent uniform training inputs $S=(X_1,\ldots,X_n)$, condition
further on $f(X_i)z_N(X_i)>0$ for every $i$.
Define the actual correct-label probability and its confidence failure by
\begin{align*}
 q_{N,S}(x)&=\mathbb P\{f(x)z_N(x)>0\mid
 |z_N(y)|>a\ \forall y;\ f(X_i)z_N(X_i)>0\ \forall i\},\\
 F_N(S)&=\mathbb P_x\{q_{N,S}(x)\le2/3\}.
\end{align*}
A successful prediction therefore requires strictly more than $2/3$
probability of the correct label. The test input is independent and
uniform.

Put $D=d+1$, $\beta=2^{-(L-1)}$, and
\begin{equation}\label{eq:main-threshold}
 B_{d,k}=\sqrt{k}\log\frac{eD}{\sqrt{k}},\qquad
 W_{d,k}=\max\{k,B_{d,k}\}.
\end{equation}
\begin{theorem}[Actual-posterior sample lower bound]\label{thm:main}
For each fixed $L\ge2$ and $a>0$, there are $c>0$, $k_0$, and $d_0$
such that, for every $d\ge d_0$, $k_0\le k\le d$, support $U$ of
size $k$, integer
\[
 0\le n\le\exp(cW_{d,k}),
\]
and fixed $0<\gamma<1$,
\[
 \mathbb P_S\left\{\liminf_{N\to\infty}F_N(S)\ge\frac1{16}\right\}
 \ge\frac34.
\]
The conditioning events have positive probability for all sufficiently
large widths at every fixed dimension and training sample.
\end{theorem}

The width limit is taken first, with $d,k,L,a,\gamma,S$ fixed.
No simultaneous width--dimension limit or quantified minimum width is
asserted. The training-sample probability holds for each specified $n$,
with constants uniform over the displayed range; no monotonicity of
fixed-target posterior risk is used.

The lower bound is exponential in the number $k$ of relevant
coordinates even when $k/d\to0$. The extra $B_{d,k}$ term is useful
when the ambient dimension is very large relative to $k$.
For $k/d$ bounded below, the theorem gives an $\exp(\Omega(d))$
sample lower bound. When $d/k\to\infty$, it does not reach
$\exp\{c k\log(ed/k)\}$. In particular it does not resolve a
$d^{\Omega(k)}$ lower bound in a polynomially sparse regime.

\subsection{Proof structure and notation}
Write $\Psi(t)=2\arcsin(t)/\pi$ and
\[
 K(x,y)=\Psi^{\circ(L-1)}\!\left(\frac{1+x\cdot y}{d+1}\right).
\]
The actual network has a Gaussian large-deviation rate with this
covariance. Conditioning on the global margin and the labels therefore
selects compatible sign functions of minimum energy
$\|z\|_K^2/2=z^{\mathsf T}K^{-1}z/2$.
We prove a simultaneous error lower bound for every minimizer, so the
argument applies to arbitrary actual weights among tied minimizers.

All function $L^p$ norms and Walsh coefficients use the uniform cube
probability measure. The characters diagonalize the normalized
covariance operator $2^{-d}K$; let $\lambda_j$ denote its degree-$j$
eigenvalue. Thus
\[
 \|z\|_K^2=\sum_{V\subseteq[d]}
             \frac{\widehat z(V)^2}{\lambda_{|V|}}.
\]
Put $D_r=\sum_{j\le r}\binom dj$, let $Q_r=K_{>r}$ retain the Fourier
degrees above $r$, and write
$\tau_r=Q_r(x,x)=\sum_{j>r}\binom dj\lambda_j$.
For an arbitrary kernel $J$, $J_{SS}$ is its sample Gram matrix.

Localized interpolation bounds every minimizing energy by
$C(d+n d^\beta)$. At the largest odd cutoff with
$n\ge A D_r\log(D_r+1)$, a one-sided squared sampling theorem controls
the low-degree part of a minimizing field.
The spectrum has equal eigenvalues at consecutive even and odd levels.
A near-orthogonal estimate controls sampled tail covariance matrices
both at the last sampled cutoff and one paired level above it. This
bounds the low-degree amplitude by a polynomial in the ratio of ambient
dimension to cutoff and makes the remaining tail small in $L^2$.
A random restriction of the active coordinates turns this head bound
into a test-error lower bound: restricting to about $(k/r)^2$ coordinates
leaves almost all of a degree-$r$ polynomial below the threshold degree
needed to approximate parity, unless its amplitude is exponentially
large in $k/r$. A direct polynomial-threshold argument supplies the
additional $B_{d,k}$ range and handles the smallest samples.

\section{The kernel rate of the actual network}
Let $M=2^d$, and identify output functions with vectors in $\mathbb R^M$.
Define
\begin{equation}\label{eq:kernel}
 \Psi(t)=\frac2\pi\arcsin t,\qquad
 K(x,y)=\Psi^{\circ(L-1)}\!\left(\frac{1+x\cdot y}{d+1}\right).
\end{equation}
The positive definiteness of $K$ is proved in the next section.
For each Boolean function $g$, put
\begin{equation}\label{eq:energy}
 E_K(g)=\min_{g(x)z(x)\ge a\ \forall x}\frac12 z^TK^{-1}z.
\end{equation}
The minimum exists since its feasible set is closed and nonempty and
the objective is coercive. Replacing the closed inequalities by strict
ones gives the same infimum, by scaling a feasible vector by $1+\epsilon$.

\begin{proposition}[Posterior concentration on minimum-energy signs]
\label{prop:rate}
For the actual network~\eqref{eq:weights}, at every fixed
$0<\gamma<1$,
\begin{equation}\label{eq:rate}
 \lim_{N\to\infty}N^{-\gamma}\log
 \mathbb P\{g(x)z_N(x)>a\ \forall x\}=-E_K(g).
\end{equation}
Consequently, after any fixed noiseless training sample, posterior
mass tends to one on the compatible functions minimizing $E_K$.
No assertion about the relative weights of tied minimizers is needed.
\end{proposition}

\begin{proof}
For hidden layer $s$, let
$Q_N^s(x,y)=N^{-1}\sum_i h_i^s(x)h_i^s(y)$.
We first establish, for each fixed $\epsilon>0$,
\begin{equation}\label{eq:freeze}
 \mathbb P\{\|Q_N^{L-1}-K\|>\epsilon\}
       \le C_\epsilon e^{-c_\epsilon N}
\end{equation}
for all sufficiently large $N$. Any matrix norm can be used, since
$M$ is fixed. Conditional on the preceding layer, the rows of the
next hidden layer are independent bounded vectors. Hoeffding's
inequality controls all their empirical second moments. Their
preactivations have covariance $N^{1-\gamma}Q_N^{s-1}$; at the first
layer the covariance is
$N^{1-\gamma}(1+x\cdot y)/d$. Because $N^{1-\gamma}\to\infty$,
the conditional second-moment map converges to the Gaussian-sign map
\[
 Q\longmapsto \left[\Psi\!\left(
 \frac{Q(x,y)}{\sqrt{Q(x,x)Q(y,y)}}\right)\right]_{x,y}.
\]
This follows by dominated convergence for a Gaussian vector, since
each scalar variance stays bounded away from zero in a neighborhood
of the limiting covariance. The convergence is uniform in such a
neighborhood: a compactness subsequence and a continuous square root
coupling reduce the assertion to the same dominated-convergence
argument. The limiting map is continuous there. Inducting through
the fixed number of layers, with smaller fixed tolerances at earlier
layers, proves~\eqref{eq:freeze}. The first-layer limit normalizes
the diagonal $(d+1)/d$, giving the kernel in~\eqref{eq:kernel}.

Conditionally on the hidden layers, the readout is exactly Gaussian:
\begin{equation}\label{eq:readout}
 Z_N\mid Q_N^{L-1}\sim\mathcal N(0,N^{-\gamma}Q_N^{L-1}).
\end{equation}
Couple it to $Y_N=N^{-\gamma/2}K^{1/2}\xi$, where
$\xi\sim\mathcal N(0,I_M)$ is independent of the hidden layers.
On a fixed sufficiently small neighborhood of the positive-definite
matrix $K$, the square-root map is Lipschitz. Hence, for $\delta>0$,
\[
 \mathbb P\{\|Z_N-Y_N\|>\delta\}
 \le C_\epsilon e^{-c_\epsilon N}
 +\mathbb P\{\|\xi\|>\delta N^{\gamma/2}/(C\epsilon)\}.
\]
First divide logarithms by $N^\gamma$ and let $N\to\infty$;
then let $\epsilon\downarrow0$. The resulting upper limit is
$-\infty$, since $\gamma<1$. The coupled vectors are therefore
exponentially equivalent at speed $N^\gamma$.

The fixed-dimensional Gaussian vectors $Y_N$ have large-deviation
rate $I(z)=z^TK^{-1}z/2$, directly from their densities. Exponential
equivalence transfers the same bounds to $Z_N$. Applying the upper
bound to a closed margin orthant and the lower bound to its interior
gives~\eqref{eq:rate}, since their infima agree. Finally, there are
only finitely many Boolean functions for fixed $d$. Dividing each
compatible orthant probability by their sum proves the posterior
concentration claim.
\end{proof}

This is a rate statement for the network itself. It does not require
relative-error approximation of rare probabilities, nor convergence
in total variation to a Gaussian posterior.


\section{Fourier level masses of the kernel}\label{app:spectrum}
Put $D=d+1$, $\beta=2^{-(L-1)}$, and $F=\Psi^{\circ(L-1)}$.
The characters $\chi_U(x)=\prod_{j\in U}x_j$ diagonalize the
normalized covariance operator $M^{-1}K$, $M=2^d$. Let $\lambda_j$
be its eigenvalue at degree $j$. We first work on the full $D$-cube
with kernel $F(u\cdot v/D)$, whose degree-$j$ eigenvalue is denoted
by $a_j$. Constants in this and the next section may depend on $L$.
Notation $b\asymp c$ means two inequalities by positive constant
multiples, uniformly in the stated ranges.

\begin{lemma}[Power coefficients and paired eigenvalues]\label{lem:power}
Write $F(z)=\sum_{m\text{ odd}}c_mz^m$. Every odd coefficient is
positive, $\sum_m c_m=1$, and
\[
 c_m\asymp\frac{\Gamma(m-\beta)}{m!}\asymp m^{-1-\beta}
 \qquad(m\ge1\text{ odd}).
\]
The full-cube eigenvalues vanish at even degrees and strictly decrease
through odd degrees. They are positive at every odd degree. On the
original $d$-cube,
\[
 \lambda_{2k}=\lambda_{2k+1}=a_{2k+1}>0,
\]
with a final unpaired eigenvalue if $d$ is even. In particular $K$
is positive definite.
\end{lemma}
\begin{proof}
The series of $\Psi$ has positive coefficients at every odd exponent,
with
\[
 [z^{2k+1}]\Psi(z)=\frac2\pi\frac{\binom{2k}{k}}{4^k(2k+1)}.
\]
Composition preserves strict positivity and oddness. Near $z=1$ in
a slit neighborhood,
\[
 \Psi(z)=1-b(1-z)^{1/2}+O((1-z)^{3/2}),
 \qquad b=\frac{2\sqrt2}{\pi}.
\]
Induction gives
$F(z)=1-A(1-z)^\beta+O((1-z)^{3\beta})$, with $A>0$;
at each composition its leading constant changes from $A_t$ to
$b\sqrt{A_t}$, starting from $A_1=b$.
These expansions admit algebraic coefficient transfer
\cite[Chapter VI]{flajolet2009}. Here are the analytic-domain details.
On the unit circle outside $\{-1,1\}$, positivity of every odd
coefficient gives $|\Psi(z)|<1$: equality in the triangle inequality
would require $z^2=1$. Hence the iterated functions extend
analytically through every compact arc away from these endpoints.
At $1$, restrict $\arg(1-z)$ to a closed proper subinterval of
$(-\pi,\pi)$. Each square-root composition reduces this angular
range, and its positive leading coefficient keeps its image in the
slit neighborhood for the next arcsine composition. Oddness gives
the identical argument at $-1$. These local sectors and the compact
arcs provide the required two-singularity indented domain.
The contributions of $1$ and $-1$ add at odd exponents, giving
\[
 c_m\sim\frac{2A}{-\Gamma(-\beta)}m^{-1-\beta}
 \quad(m\to\infty\text{ odd}).
\]
Together with positivity of the finitely many initial coefficients
and the gamma-ratio asymptotic, this proves the coefficient comparisons.
For $L=2$ the explicit formula and Stirling's bounds suffice instead.

Let $p_m(j)$ be the probability that the simple coordinate-flip walk
on the $D$-cube reaches a specified distance-$j$ vertex in $m$ steps.
Expanding $(u\cdot v/D)^m$ shows that its degree-$j$ Fourier coefficient
is $p_m(j)$. Thus $a_j=\sum_{m\text{ odd}}c_mp_m(j)$.
For the walk recursion,
\begin{align*}
 D[p_{m+1}(j)-p_{m+1}(j+2)]
 ={}&j[p_m(j-1)-p_m(j+1)]\\
 &+(D-j-2)[p_m(j+1)-p_m(j+3)].
\end{align*}
Induction from time zero gives nonincreasing probabilities through
each parity class; terms with zero boundary coefficients are omitted.
The term $m=j$ makes the odd-eigenvalue inequalities strict and proves
positivity. Even-degree eigenvalues vanish by parity.
Restricting the distinguished coordinate to $+1$ gives
$\lambda_q=a_q+a_{q+1}$, proving the claims.
\end{proof}

For odd $r$, the exact combined multiplicity of reduced degrees
$j-1,j$ is $\binom d{j-1}+\binom dj=\binom Dj$. Consequently
\begin{equation}\label{eq:paired}
 D_r=\sum_{\substack{j\le r\\j\text{ odd}}}\binom Dj,
 \quad
 \tau_r=\sum_{\substack{j>r\\j\text{ odd}}}\binom Dj a_j,
 \quad \lambda_{r+1}=a_{r+2}.
\end{equation}
The odd cutoff is essential for these exact identities.

\begin{lemma}[Sharp level and tail estimates]\label{lem:sharp}
Uniformly for odd $1\le j\le D/4$ and odd $1\le r\le D/8$,
\begin{align}
 \binom Dj a_j&\asymp j^{-1-\beta},\label{eq:level}\\
 \tau_r&\asymp r^{-\beta},\label{eq:trace}\\
 \frac{\lambda_{r+1}}{\tau_r}
 &\asymp\frac1{r\binom D{r+2}},\label{eq:ratio}\\
 s_{2,r}:=\sum_{q>r}\binom dq\lambda_q^2
 &\le\frac{C r^{-2\beta-2}}{\binom D{r+2}}.
 \label{eq:secondmoment}
\end{align}
The same level-mass estimate holds up to $j\le\rho D$ for any fixed
$\rho<1/2$, with constants depending on $\rho$.
\end{lemma}
\begin{proof}
The gamma integral and Lemma~\ref{lem:power} give, for odd $j$,
\begin{align}
 a_j&\asymp D^{-\beta}\int_0^\infty
 u^{-1-\beta}p(u)^j(1-p(u))^{D-j}\,du,
 \quad p(u)=\frac{1-e^{-2u}}2,\label{eq:poisson}\\
 &\asymp D^{-\beta}\int_0^{1/2}p^j(1-p)^{D-j}w_\beta(p)\,dp,
 \label{eq:betaweight}
\end{align}
where $w_\beta(p)=[-\frac12\log(1-2p)]^{-1-\beta}/(1-2p)$.
Indeed the gamma integral Poissonizes the walk time, and independent
Poisson coordinate counts give independent coordinate parities with
success probability $p(u)$. Since an odd-distance vertex has zero
probability at even walk times, no extra parity factor occurs.
The weight $w_\beta$ is integrable near $1/2$ and comparable to
$p^{-1-\beta}$ on $[0,b]$ for every fixed $b<1/2$.

Fix $\rho<b<1/2$ and take $j\le\rho D$. The integral on $[0,b]$
is comparable to the complete beta integral
$B(j-\beta,D-j+1)$. For the lower comparison, a beta random variable
with these parameters has mean at most $\rho$, so Markov's inequality
puts at least a fraction $1-\rho/b$ of its mass in $[0,b]$.
The identity
\[
 \binom Dj B(j-\beta,D-j+1)
 =\frac{\Gamma(D+1)\Gamma(j-\beta)}
        {\Gamma(j+1)\Gamma(D+1-\beta)}
 \asymp D^\beta j^{-1-\beta}
\]
gives the desired mass. On $[b,1/2]$, the binomial factor decreases
with $p$ and, after multiplication by $\binom Dj$, is at most
$\exp[-D\operatorname{KL}(\rho\Vert b)]$. The remaining integrable
weight makes this contribution exponentially small. This proves
\eqref{eq:level} and its stated extended range.

Sum \eqref{eq:poisson} over the odd distances above $r$. Dropping the
parity restriction and using the binomial Markov bound gives
\[
 \tau_r\le C D^{-\beta}\int_0^\infty
 u^{-1-\beta}\min\{1,Du/(r+1)\}\,du
 \le C(r+1)^{-\beta}.
\]
For the lower bound, sum \eqref{eq:level} over the odd $j$ in
$(r,2r+2]$, using a fixed cutoff strictly below $D/2$ as needed at
the endpoint. This yields $c r^{-\beta}$; bounded $r$ are absorbed
by a change of constants. Equation~\eqref{eq:ratio} follows from
\eqref{eq:level}, \eqref{eq:trace}, and \eqref{eq:paired}.

For the second moment, use paired multiplicities to write
$s_{2,r}=\sum_{j>r,\ j\text{ odd}}\binom Dj a_j^2$.
Up to an odd cutoff $b_D$ near $D/4$, its terms are at most
$Cj^{-2\beta-2}/\binom Dj$. Consecutive odd inverse binomial
coefficients have ratio
\[
 \frac{(j+1)(j+2)}{(D-j)(D-j-1)},
\]
uniformly less than a fixed constant below one in this range.
The sum is therefore bounded by a constant times its first term,
at $j=r+2$. Beyond $b_D$, monotonicity gives the bound
$a_{b_D}\tau_{b_D}\le C D^{-2\beta-1}/\binom D{b_D}$,
which is exponentially smaller than the claimed bound uniformly
for $r\le D/8$. This proves \eqref{eq:secondmoment}.
\end{proof}

Finally, two-step binomial ratios give, uniformly for odd $r\le D/8$,
\begin{equation}\label{eq:dimensionratios}
 \binom Dr\le D_r\le C\binom Dr,
 \qquad \frac{D_r}{\binom D{r+2}}\le C(r/D)^2.
\end{equation}
All constants used here can be fixed independently of a subsequently
chosen small positive limit $r/d\to\alpha$.

\begin{corollary}[Lower bounds used by the interpolation]\label{cor:lower-spectrum}
For sufficiently large $d$,
\[
 \lambda_0\ge c/d,\qquad
 \lambda_j\ge \frac{c d^{-1-\beta}}{\binom dj}\quad(0\le j\le .49d),
 \qquad v:=\lambda_{\min}(K)\ge c(d\log d)^{-\beta}.
\]
Also $v\to0$.
\end{corollary}
\begin{proof}
The first bound follows from the level estimate at full-cube degree one.
For odd reduced degree $j$, use $\lambda_j=a_j$ and
$\binom Dj/\binom dj=D/(D-j)$. For even $j$, use
$\lambda_j=a_{j+1}$ and
$\binom D{j+1}/\binom dj=D/(j+1)$.
The extended level bound with any fixed cutoff between $.49$ and $1/2$
then gives the stated uniform lower bound, including $j=0$.

For any odd full-cube degree $j\le D$, restrict the integral
\eqref{eq:poisson} to $u\ge\tfrac12\log(2D)$. In this range
$p(u)\ge(1-1/(2D))/2$ and $1-p(u)\ge1/2$, so
$p(u)^j(1-p(u))^{D-j}\ge c2^{-D}$. Integrating $u^{-1-\beta}$ gives
$a_j\ge c2^{-D}D^{-\beta}(\log D)^{-\beta}$.
Every reduced eigenvalue equals an odd $a_j$, and the matrix eigenvalues
are $M$ times the normalized ones. Thus $v\ge c(d\log d)^{-\beta}$.
Finally, for any fixed small positive cutoff fraction, the minimum
normalized eigenvalue is at most the average in its tail, whence
$v\le M\tau_r/(M-D_r)=O(d^{-\beta})$. This proves $v\to0$.
\end{proof}

\section{A globally separated interpolation comparator}
\label{sec:comparator}
The next construction bounds the energy of every minimizer. It works
for arbitrary training labels and enforces the margin on the entire cube.

\begin{lemma}[Localized interpolation]\label{lem:comparator}
For $1\le n\le e^{d/25}$, with probability at least $1-e^{-c_Ld}$
over independent uniform inputs, the following holds simultaneously for
every labeling $y_i\in\{-1,1\}$ and every $A>0$.
There is a field $h$ with
\[
 |h(x)|>A\quad\text{for every }x,\qquad
 y_i h(X_i)>A\quad\text{for every }i,
 \qquad \|h\|_K^2\le C_L A^2(d+n d^\beta).
\]
For $n=0$, the same deterministic energy bound holds using a constant field.
\end{lemma}
\begin{proof}
For $z\in\mathcal X$ define the nonnegative bump
\[
 b_z(x)=20\int_{2/5}^{9/20}(1-2p)^{\operatorname{dist}(x,z)}\,dp.
\]
It satisfies $b_z(z)=1$ and $b_z(x)\le(1/5)^{\operatorname{dist}(x,z)}$.
Its degree-$j$ Walsh coefficients, apart from their translation signs,
are
\[
 \widehat b_j=20\int_{2/5}^{9/20}(1-p)^{d-j}p^j\,dp,
 \qquad
 \binom dj\widehat b_j\le\frac{20}{d+1},\qquad
 \widehat b_j\le(3/5)^d.
\]
The middle inequality follows by extending to the complete beta integral.
By Corollary~\ref{cor:lower-spectrum},
$\lambda_j\ge c_Ld^{-1-\beta}/\binom dj$ for $j\le.49d$.
Hence these degrees contribute at most
\[
 C_Ld^{1+\beta}\sum_{j\le.49d}
       \left[\binom dj\widehat b_j\right]^2\le C_Ld^\beta
\]
to $\|b_z\|_K^2$.
For $j>.49d$, the integrand is increasing in $p$ on the integration
interval, so $\widehat b_j\le.55^{d-j}.45^j$.
Let $v=\lambda_{\min}(K)$ be the smallest matrix eigenvalue.
The same corollary gives $v\ge c_L(d\log d)^{-\beta}$ and
$\lambda_j\ge v/2^d$. Thus the remaining energy is at most
\[
 \frac{2^d}{v}\sum_{j>.49d}\binom dj .55^{2(d-j)}.45^{2j}
 \le\operatorname{poly}(d)e^{-c_*d},
\]
where
\[
 c_*=-[\log2+\mathsf h(.49)+1.02\log(.55)+.98\log(.45)]>.0062,
 \qquad \mathsf h(t)=-t\log t-(1-t)\log(1-t).
\]
The exponent decreases beyond $.49d$, so the binomial entropy bound
applies. We have proved $A_d:=\|b_z\|_K^2\le C_Ld^\beta$.

Put $J(z,w)=\langle b_z,b_w\rangle_K$.
Its normalized Fourier eigenvalues are
$\nu_j=\widehat b_j^2/\lambda_j\le .72^d/v$.
Parseval and Cauchy--Schwarz give
\[
 \mathbb E_{X,Y}|J(X,Y)|
 \le\left(\sum_j\binom dj\nu_j^2\right)^{1/2}
 \le\sqrt{\nu_{\max}A_d}.
\]
For every subset $T$ of the sampled indices, including one chosen after
the sample is seen,
\[
 \left\|\sum_{i\in T}b_{X_i}\right\|_K^2
 \le nA_d+\sum_{i\ne j}|J(X_i,X_j)|.
\]
The cross sum has expectation at most
$n(n-1)\sqrt{\nu_{\max}A_d}$.
Markov's inequality at threshold $n d^\beta$ has exponentially small
failure probability for $n\le e^{.04d}$, because
$\tfrac12\log(1/.72)>.16$.
Thus all subsets simultaneously have energy at most $C_Ln d^\beta$.

Also, except with exponentially small probability, every pair of
different sample indices has inputs at Hamming distance greater than
$.2d$. A union bound costs at most
\[
 n^2\exp[-d(\log2-\mathsf h(.2))],
 \qquad \log2-\mathsf h(.2)>.19.
\]
This event excludes repeated inputs.
At any input at most one center is within distance $.1d$; the sum of
all other bumps is at most $n(.2)^{.1d}=o(1)$.
Consequently, for every selected subset of centers, its bump sum is
$1+o(1)$ at a selected center and at most $.2+o(1)$ everywhere else.

For $T=\{i:y_i=-1\}$, take
\[
 h=2A\left(1-2\sum_{i\in T}b_{X_i}\right).
\]
The preceding estimates prove its strict global and training margins
for all sufficiently large $d$.
Finally $\|1\|_K^2=1/\lambda_0\le C_Ld$ by
Corollary~\ref{cor:lower-spectrum}; the squared triangle inequality
proves the asserted energy bound. For $n=0$, use $h=2A$.
\end{proof}

\section{One-sided squared sampling with one logarithm}
\label{sec:sampling}

The sampling estimate needed below concerns a squared one-sided loss.
It follows from a metric estimate in the proof of Lee's chaining theorem
\cite[Lemma~2.12, equation~(2.20)]{Lee2023}. We spell out the comparison:
a first-power ReLU sampling theorem alone would not give this result.
Throughout this section, $C$ denotes a universal constant that may change
from line to line, and $t_-:=\max\{-t,0\}$.

\begin{lemma}[A squared Lipschitz process]\label{lem:squared-process}
Let $a_1,\ldots,a_m\in\mathbb R^D$, and set
\[
 R=\max_{i\le m}\|a_i\|_2,
 \qquad V=\sup_{\|u\|_2\le1}
       \left(\sum_{i=1}^m\langle a_i,u\rangle^2\right)^{1/2}.
\]
For each $i$, let $\psi_i:\mathbb R\to\mathbb R$ satisfy
$\psi_i(0)=0$ and $\operatorname{Lip}(\psi_i)\le1$.
For independent uniform signs $\epsilon_i$,
\begin{equation}\label{eq:squared-process}
 \mathbb E_\epsilon\sup_{\|u\|_2\le1}
 \left|\sum_{i=1}^m\epsilon_i\psi_i(\langle a_i,u\rangle)^2\right|
 \le C R\sqrt{\log(2(m+D))}\,V.
\end{equation}
\end{lemma}

\begin{proof}
Use the genuine seminorms $N_i(u)=|\langle a_i,u\rangle|$ and the
symmetric quasimetric
\[
 d_A(u,v)^2
 =\sum_{i=1}^m N_i(u-v)^2\bigl(N_i(u)^2+N_i(v)^2\bigr).
\]
The stronger intermediate estimate established in the proof of
\cite[Lemma~2.12]{Lee2023} is
\begin{equation}\label{eq:lee-metric}
 \gamma_2(T,d_A)
 \le C\bigl(\lambda\sqrt{\log(2D)}+\kappa\bigr)
       \sup_{u\in T}\left(\sum_i N_i(u)^2\right)^{1/2},
 \qquad T\subseteq B_2^D,
\end{equation}
where
\[
 \lambda=\max_i\bigl(\mathbb E N_i(g)^2\bigr)^{1/2},
 \qquad \kappa=\mathbb E\max_i N_i(g),
 \qquad g\sim\mathcal N(0,I_D).
\]
Here $\gamma_2$ is the generic-chaining functional. This is equation
(2.20) of that proof, with the energy normalization (2.17) undone.
Seminorms follow from the norm statement by replacing
$N_i(u)$ with $(N_i(u)^2+\varepsilon^2\|u\|_2^2)^{1/2}$ and then
letting $\varepsilon\downarrow0$; the resulting metrics dominate
$d_A$, and the parameters on the right converge.
For our rank-one seminorms,
$\lambda=R$ and $\kappa\le R\sqrt{2\log(2m)}$.

Let $X_u=\sum_i\epsilon_i\psi_i(\langle a_i,u\rangle)^2$.
Its subgaussian increment metric is bounded by $\sqrt2\,d_A$, because
\begin{align*}
 &\sum_i\left(\psi_i(\langle a_i,u\rangle)^2
             -\psi_i(\langle a_i,v\rangle)^2\right)^2\\
 &\quad\le\sum_i |\langle a_i,u-v\rangle|^2
       \bigl(|\langle a_i,u\rangle|+|\langle a_i,v\rangle|\bigr)^2
 \le 2d_A(u,v)^2.
\end{align*}
We used both Lipschitz continuity and
$|\psi_i(t)|\le|t|$. The subgaussian chaining inequality, together
with $X_0=0$, bounds $\mathbb E\sup X_u$ by
$C\gamma_2(B_2^D,d_A)$. Apply the same argument to $-X_u$ and add.
Equation~\eqref{eq:lee-metric} now proves
\eqref{eq:squared-process}.
\end{proof}

\begin{remark}[The normalization in the metric estimate]
The proof of \cite[Lemma~2.12]{Lee2023} passes through a coarse additive
constant in its equation (2.26), after normalizing
$\sup_T\sum_iN_i^2=1$. The homogeneous version
\eqref{eq:lee-metric} follows by retaining the scale in the two geometric
terms of that argument. Indeed, in the rank-one case $T\subseteq B_2^D$
implies $N_i(u)\le\lambda$ on $T$; writing
$\|u\|_{\mathcal N}=\max_iN_i(u)$ gives
$d_A(u,v)\le\sqrt2\|u-v\|_{\mathcal N}$ on $T$.
Thus $\operatorname{diam}_{d_A}(T)\le C\lambda$ and the volumetric
packing constant in its Corollary~2.11 can be taken to be
$C\lambda$, since $T\subseteq\lambda B_{\mathcal N}$.
The additive term is therefore $C\lambda$, absorbed by
$C(\lambda\sqrt{\log(2D)}+\kappa)$.
The growth-functional portion of the proof is unchanged.
If the packing constant is less than one, first rescale the quasimetric
to apply that corollary in its stated normalization, and then scale back.
Most importantly, throughout that argument the $N_i$ are symmetric
seminorms. No assertion that a positive part is a seminorm is used.
\end{remark}

\begin{lemma}[Uniform one-sided squared sampling]
\label{lem:squared-sampling}
Let $(Z,Y)$ be a random pair with $Z\in\mathbb R^D$, $Y\in\{-1,1\}$,
$\mathbb E ZZ^{\mathsf T}=I_D$, and $\|Z\|_2^2=D$ almost surely.
For independent copies $(Z_i,Y_i)$, let
\[
 \Delta_n=\sup_{\|u\|_2\le1}
 \left|\frac1n\sum_{i=1}^n(Y_i\langle Z_i,u\rangle)_-^2
          -\mathbb E(Y\langle Z,u\rangle)_-^2\right|,
 \qquad h_n=\sqrt{\frac{D\log(2(n+D))}{n}}.
\]
Then $\mathbb E\Delta_n\le C h_n(1+C h_n)$.
In particular, for any fixed $\varepsilon,\delta>0$ there is a finite
$A_{\varepsilon,\delta}$ such that
\[
 n\ge A_{\varepsilon,\delta}D\log(2D)
 \quad\Longrightarrow\quad
 \mathbb P\{\Delta_n>\varepsilon\}\le\delta.
\]
The labels may depend arbitrarily on the features.
\end{lemma}

\begin{proof}
Put $Q_n=n^{-1}\sum_iZ_iZ_i^{\mathsf T}$.
Conditional on the observations, apply
Lemma~\ref{lem:squared-process} to
$a_i=Y_iZ_i/\sqrt n$ and $\psi_i(t)=t_-$.
Here $R=\sqrt{D/n}$ and $V=\sqrt{\|Q_n\|_{\mathrm{op}}}$.
Symmetrization gives
\begin{equation}\label{eq:sampling-symmetrization}
 \mathbb E\Delta_n
 \le C h_n\,\mathbb E\sqrt{\|Q_n\|_{\mathrm{op}}}.
\end{equation}
We also apply the same process lemma with $\psi_i(t)=t$.
Since $\mathbb E Q_n=I_D$, symmetrization and Jensen's inequality imply,
with $A=\mathbb E\|Q_n\|_{\mathrm{op}}$,
\[
 A\le1+\mathbb E\|Q_n-I_D\|_{\mathrm{op}}
 \le1+C h_n\mathbb E\sqrt{\|Q_n\|_{\mathrm{op}}}
 \le1+C h_n\sqrt A.
\]
Solving this quadratic inequality yields $\sqrt A\le1+C h_n$,
after increasing $C$. Substitute into
\eqref{eq:sampling-symmetrization}.
Finally, at $n\ge A_{\varepsilon,\delta}D\log(2D)$ the quantity
$h_n$ can be made uniformly small by increasing
$A_{\varepsilon,\delta}$. Markov's inequality proves the probability
statement. It is a statement for each specified sample size $n$.
\end{proof}


\section{A polynomial-threshold obstruction}
The following elementary polynomial-threshold obstruction is due to
Aspnes, Beigel, Furst, and Rudich~\cite{ABFR}; we include its proof.
\begin{lemma}[Polynomial thresholds cannot approximate parity]
\label{lem:ptf-parity}
On a cube of dimension $k$, for $0\le s<k$, every degree-$s$ polynomial threshold, with a fixed common convention
for zero ties, has full-parity error at least
\begin{equation}\label{eq:ptf-parity}
 2^{-k}\sum_{j=0}^{\lfloor(k-s-1)/2\rfloor}\binom kj.
\end{equation}
This lower bound tends to $1/2$ whenever $s=o(\sqrt k)$.
\end{lemma}
\begin{proof}
A sufficiently small constant perturbation resolves any zero ties
without changing the specified threshold classifier. Thus suppose
$P$ has no zeros. Put $t=\lfloor(k-s-1)/2\rfloor$. If its error set
had fewer than $\sum_{j\le t}\binom kj$ points, a nonzero polynomial
$Q$ of degree at most $t$ could be chosen to vanish on every error.
The function $fPQ^2$ would then be nonnegative everywhere and positive
somewhere. But $PQ^2$, after multilinear reduction, has degree at most
$s+2t<k$, so its parity Fourier coefficient
$\mathbb E fPQ^2$ is zero, a contradiction. This proves
\eqref{eq:ptf-parity}. The last assertion follows from the central
binomial bound $\max_j2^{-k}\binom kj=O(k^{-1/2})$.
\end{proof}


\section{A degree--amplitude inequality from random restrictions}
\label{sec:random-restriction}
The next lemma uses the global margin to strengthen a degree-only
obstruction. Its proof combines random restrictions with the parity
approximation bound of Aspnes, Beigel, Furst, and Rudich.

\begin{lemma}[Robust degree--amplitude inequality]
\label{cor:restriction-amplitude}
\label{lem:random-restriction-exact}
Let $f=\chi_U$ on the uniform $d$-cube, where $|U|=k$, and let
$|z(x)|\ge a>0$ everywhere. Suppose $z=P+T$, where $\deg P\le u$,
$1\le u\le k$, and $P,T$ have disjoint Fourier supports. There are
absolute constants $c,C>0$ such that
\begin{equation}
 \Pr\{\operatorname{sign}z\ne f\}
 \ge \frac7{16}-\frac{\|T\|_2^2}{a^2}
       -C\frac{\|P\|_2^2}{a^2}\exp\!\left(-c\frac{k}{u}\right).
 \label{eq:restriction-amplitude}
\end{equation}
More precisely, for integers $1\le m\le k$ and $0\le q<m$ the right-hand
side can be replaced by
\begin{equation}
 2^{-m}\sum_{j=0}^{\lfloor(m-q-1)/2\rfloor}\binom mj
 -\frac{\|T\|_2^2}{a^2}
 -\frac{\|P\|_2^2}{a^2}
    \min\!\left\{1,
       \left(\frac{eum}{k(q+1)}\right)^{q+1}\right\}.
 \label{eq:restriction-exact}
\end{equation}
\end{lemma}
\begin{proof}
Choose a uniform $m$-element subset $J\subseteq U$, fix the coordinates
outside $J$ uniformly, and truncate the restricted polynomial $P$ to
degree $q$ in the remaining coordinates. Denote the resulting polynomial
by $Q$. Viewed on the original cube, this operation retains exactly the
Fourier terms $\widehat P(S)\chi_S$ with $|S\cap J|\le q$. Orthogonality
therefore gives
\[
 \mathbb E_J\|z-Q\|_2^2
 =\|T\|_2^2+
   \sum_S\widehat P(S)^2\Pr_J\{|S\cap J|>q\}.
\]
Since $|S|\le u$, the probability is bounded by the upper tail of a
hypergeometric random variable with population size $k$, at most $u$
marked elements, and sample size $m$. The factorial-moment estimate gives
\[
 \Pr\{|S\cap J|\ge q+1\}
 \le \min\!\left\{1,
       \left(\frac{eum}{k(q+1)}\right)^{q+1}\right\}.
\]
For every restriction, the target is parity or its negation on the
$m$ remaining coordinates. Lemma~\ref{lem:ptf-parity} gives error
at least
\[
 B(m,q):=2^{-m}\sum_{j=0}^{\lfloor(m-q-1)/2\rfloor}\binom mj
\]
for every degree-$q$ polynomial threshold function on that cube. This
statement holds for real coefficients: its proof uses real linear
algebra. Ties can be assigned either fixed sign by adding an arbitrarily
small constant. Whenever $z$ and $Q$ have different signs, the global
margin of $z$ implies $|z-Q|\ge a$. Averaging the resulting error bound
proves~\eqref{eq:restriction-exact}.

For completeness, choose a sufficiently small absolute $\delta>0$,
put $t=k/u$, and, when $t$ exceeds a sufficiently large absolute constant,
set $m=\min\{k,\lfloor\delta t^2\rfloor\}$. If $m=k$, then
$u\le\sqrt{\delta k}$; use $Q=P$ directly. The central binomial
atom bound yields $B(k,u)\ge7/16$ after increasing the fixed lower
threshold on $t$. If $m<k$, take $q=\lfloor\sqrt m/32\rfloor$.
For example $\delta=(1024e)^{-2}$ ensures
\[
 \frac{eum}{k(q+1)}\le32e\sqrt\delta\le\frac1{32},
 \qquad q+1\ge c_0 t,
 \qquad B(m,q)\ge\frac7{16}
\]
for sufficiently large $t$, with $c_0>0$ absolute. This proves
\eqref{eq:restriction-amplitude} in that range. For bounded $t$,
use $\|P\|_2^2+\|T\|_2^2=\|z\|_2^2\ge a^2$ and enlarge $C$ so
that the asserted lower bound is nonpositive.
\end{proof}

In particular, any degree-$u$ polynomial with global margin $a$ whose
threshold has parity error at most $3/8$ must satisfy
\[
 \|P\|_2^2\ge c' a^2\exp(c k/u).
\]
For a fixed degree fraction $u/k$, this lower bound is a constant. It is
therefore consistent with the bounded-amplitude blocked construction in
Proposition~\ref{prop:blocked-approximants}.

% Primary reference to add to the bibliography:
% James Aspnes, Richard Beigel, Merrick Furst, and Steven Rudich.
% The expressive power of voting polynomials. Combinatorica 14(2), 1994.
% Theorem 2.2 and Lemma 2.3.
% https://www.cs.yale.edu/homes/aspnes/papers/stoc91voting.pdf

\section{Random tail covariance matrices}\label{app:gram}
Let $Q$ be the covariance kernel retaining only degrees greater than an
odd cutoff $r$, and write $\tau=\tau_r$, $\ell=\lambda_{r+1}$, and
$s_2=s_{2,r}$. This section proves the sample event needed later, without
assuming that a Gaussian approximation remains valid after conditioning.
The entire argument concerns the exact limiting Gaussian covariance.

\begin{lemma}[Far-pair decay]\label{lem:far}
There are constants $c_0,C>0$ such that, when $r\le D/8$ is odd,
\[
 \max_{d/4\le\dist(x,y)\le3d/4}|Q(x,y)|
 \le\rho_r:=C\sqrt D\,e^{-r/8}+Ce^{-c_0D}.
\]
\end{lemma}
\begin{proof}
The full-cube Krawtchouk polynomial satisfies
\[
 \sum_{j=0}^D\mathsf K_j(h;D)z^j=(1+z)^{D-h}(1-z)^h.
\]
If $\delta D\le h\le(1-\delta)D$, put $b=\min(h,D-h)$.
On $|z|=t$ the generating function is bounded in absolute value by
\[
 (1+t^2)^b(1+t)^{D-2b}
 \le(1+t)^D\exp\!\left[-\frac{2\delta Dt}{(1+t)^2}\right].
\]
Cauchy's coefficient bound at $t=j/(D-j)$ and the uniform Stirling
lower bound for $\binom Dj$ therefore give, for $1\le j<D$,
\begin{equation}\label{eq:kraw}
 \frac{|\mathsf K_j(h;D)|}{\binom Dj}
 \le C\sqrt D\exp[-2\delta j(1-j/D)].
\end{equation}
For reduced-cube points at distance in $[d/4,3d/4]$, their augmented
full-cube distance lies in $[D/8,7D/8]$ for large $d$.
Because $r$ is odd, the exact tail is
$Q(x,y)=\sum_{j>r,\ j\text{ odd}}a_j\mathsf K_j(h;D)$.
For $r<j\le D-r$, \eqref{eq:kraw} gives a bound
$C\sqrt D e^{-r/8}$ after summation, since the total spectral mass
is one. For $j>D-r$, monotonicity and an odd cutoff $b_D$ near $D/4$
give
\[
 \sum_{j>D-r}\binom Dj a_j
 \le a_{b_D}\sum_{k<r}\binom Dk
 \le C D^{-\beta-1/2}
       \exp[-D(\mathsf h(1/4)-\mathsf h(1/8))+O(1)].
\]
Here $\mathsf h(t)=-t\log t-(1-t)\log(1-t)$ denotes the binary entropy in nats.
The positive entropy gap absorbs the polynomial factor and proves the
claim. Adjust constants for bounded $d$.
\end{proof}

\begin{lemma}[A tail Gram deviation bound]\label{lem:gram-general}
Let $X_1,\ldots,X_n$ be independent uniform cube points, with
$1\ll n\le M/4$. Let
\[
 B_{\mathrm{near}}=\#\{y:\dist(x,y)\notin[d/4,3d/4]\}
 \le 2e^{d\mathsf h(1/4)},\qquad q=2\log M.
\]
The count is independent of $x$. For every $z>0$,
\begin{align}
 &\mathbb P\{\|Q_{SS}-\tau I_n\|_{\mathrm{op}}>z\tau\}
 \notag\\
 &\quad\le\frac C{z^2}
 \left[\frac{q\rho_r}{\tau}
       +\frac{\sqrt{nq s_2}}{\tau}
       +\frac{n\ell}{\tau}
       +\frac{n(1+B_{\mathrm{near}})}M\right]^2
       +\frac{n^2(1+B_{\mathrm{near}})}M.
 \label{eq:gram-bound}
\end{align}
Here the sampled matrix includes repeated training locations.
\end{lemma}
\begin{proof}
Let $J=Q-\tau I_M$, a zero-diagonal Hermitian matrix. Form a deterministic
mask $J_{\mathrm{far}}$ retaining only pairs at distances in
$[d/4,3d/4]$. Since $Q$ is positive semidefinite with constant diagonal
$\tau$, every entry has absolute value at most $\tau$. Thus
\begin{align*}
 \|J_{\mathrm{far}}\|_{\max}&\le\rho_r,\\
 \|J_{\mathrm{far}}\|_{1\to2}^2&\le M s_2,\\
 \|J_{\mathrm{far}}\|_{\mathrm{op}}
 &\le M\ell+\tau+\tau B_{\mathrm{near}}.
\end{align*}
For the second inequality, Parseval gives squared column norm
$M s_2-\tau^2$ for $J$, and entrywise masking only decreases it.
For the third, $\|J\|\le M\ell+\tau$, while the removed matrix has
maximum absolute row sum at most $\tau B_{\mathrm{near}}$.
The entropy bound on $B_{\mathrm{near}}$ is the elementary binomial
tail bound on distances to a fixed cube point.

Let $R$ retain each of the $M$ coordinates independently with
probability $p=2n/M$. The zero-diagonal case of Tropp's random
principal-submatrix theorem \cite[Theorem 1.1]{tropp2008}, at moment
order two and $q=2\log M$, states
\[
 (\mathbb E\|RJR\|^2)^{1/2}
 \le C\left[q(\mathbb E\|RJR\|_{\max}^2)^{1/2}
       +\sqrt{pq}(\mathbb E\|JR\|_{1\to2}^2)^{1/2}
       +p\|J\|\right].
\]
Apply it to $J_{\mathrm{far}}$ and substitute the preceding bounds:
\begin{equation}\label{eq:bernoulli-gram}
 (\mathbb E\|R J_{\mathrm{far}}R\|^2)^{1/2}
 \le C[q\rho_r+\sqrt{nq s_2}+n\ell
               +(n/M)\tau(1+B_{\mathrm{near}})].
\end{equation}
The logarithm here is that of the original dimension $M$, not the
sample size.

To pass to a uniform subset of size $n$, use a uniform permutation of
all cube points, independently of $B\sim\operatorname{Bin}(M,2n/M)$.
Its first $B$ elements form the Bernoulli subset. On $B\ge n$, the
first $n$ principal submatrix is a compression of the larger one, so
its operator norm cannot be larger. The size-$n$ event is independent
of $B$, and $\mathbb P(B\ge n)\ge1/2$ for large $n$. Thus Markov's
inequality applied to \eqref{eq:bernoulli-gram} gives the same fixed-size
probability bound with a constant factor loss.

Conditional on no collisions, an independent size-$n$ sample is a
uniform size-$n$ subset, up to an irrelevant ordering. Collisions have
probability at most $n(n-1)/(2M)$. The probability of any sampled pair
violating the distance condition is at most $n^2B_{\mathrm{near}}/M$.
On their complement, the sampled matrices $J$ and $J_{\mathrm{far}}$
agree. Adding these exceptional probabilities proves
\eqref{eq:gram-bound}.
\end{proof}


\section{Tail covariance estimates at sparse cutoffs}\label{sec:sparse-gram}
The original principal-submatrix estimate uses the logarithm of the full
cube cardinality. We retain that estimate here. A more precise bound on
entries near orthogonal input pairs makes it applicable even when the
spectral cutoff is a sufficiently large fixed integer.

\begin{lemma}[Near-orthogonal entries]\label{lem:near-orthogonal}
There are constants $a_0,\delta_0,c,C>0$ such that the following holds.
Put $D=d+1$, let $r$ be odd with $r+2\le\delta_0D$, and let $h$ be the
Hamming distance between the augmented inputs $(1,x),(1,y)$.
If $a=|D-2h|/D\le a_0$, then
\begin{equation}\label{eq:near-orthogonal}
 \frac{|Q_r(x,y)|}{\tau_r}
 \le [C(a+\sqrt{(r+2)/D})]^{r+2}+C d^\beta e^{-cd},
 \qquad Q_r=K_{>r}.
\end{equation}
\end{lemma}
\begin{proof}
For $1\le j<D$, the Krawtchouk generating function gives
\begin{equation}\label{eq:sparse-kraw}
 \frac{|\mathsf K_j(h;D)|}{\binom Dj}
 \le [e^{3/2}(a+\sqrt{j/D})]^j.
\end{equation}
Indeed, set $b=\min(h,D-h)$. On the circle $|z|=t$, the generating
function has modulus at most
\[
 (1+t^2)^b(1+t)^{D-2b}\le \exp(Dt^2/2+Dat).
\]
Cauchy's coefficient bound and $\binom Dj\ge(D/j)^j$, with
$t=j/[D(a+\sqrt{j/D})]$, prove \eqref{eq:sparse-kraw}: both $Dt^2$
and $Dat$ are at most $j$.

Choose $a_0,\delta_0$ so small that
$e^{3/2}(a_0+\sqrt{\delta_0})\le e^{-1}$ and $\delta_0<1/8$.
On $j\le\delta_0D$, the derivative of
$j\log[e^{3/2}(a+\sqrt{j/D})]$ is at most $-1+1/2<0$.
Consequently \eqref{eq:sparse-kraw} is bounded by its value at $r+2$
throughout $r<j\le\delta_0D$. The sum of the corresponding positive
spectral masses $\binom Dj a_j$ is at most $\tau_r$.
For $\delta_0D<j\le(1-\delta_0)D$, the central-distance estimate
\eqref{eq:kraw} bounds the normalized Krawtchouk polynomial by
$C\sqrt D e^{-cD}$; its polynomial factor can be absorbed by decreasing
$c$. For $j>(1-\delta_0)D$, choose an odd $b_D$ near $D/4$.
Monotonicity of $a_j$ and Lemma~\ref{lem:sharp} bound the remaining
mass by
\[
 a_{b_D}\sum_{i<\delta_0D}\binom Di
 \le \operatorname{poly}(D)
       e^{-D[\mathsf h(1/4)-\mathsf h(\delta_0)]},
\]
where $\mathsf h(t)=-t\log t-(1-t)\log(1-t)$.
The entropy gap is positive. Finally $\tau_r\ge cD^{-\beta}$ by
Lemma~\ref{lem:sharp}, proving \eqref{eq:near-orthogonal}.
\end{proof}

\begin{lemma}[Sparse adaptive tail Gram matrices]\label{lem:sparse-gram}
Fix $A\ge1$, a failure probability $\delta>0$, and an odd integer
$r_0\ge11$. There is $\alpha_0>0$ such that, for any fixed
$0<\alpha\le\alpha_0$, the following assertions hold uniformly as
$d\to\infty$. Suppose $r$ is odd,
\[
 r_0\le r\le\alpha d,\qquad
 A D_r\log(D_r+1)\le n<A D_{r+2}\log(D_{r+2}+1).
\]
Put $u=r+2$, $\rho=r/(d+1)$, and $H=\log(e/\rho)$.
For each such specified $n$, with probability at least $1-\delta$,
\begin{align}
 \|(Q_r)_{SS}\|_{\mathrm{op}}
 &\le C_{L,\delta}\left[1+\frac{\sqrt{nd\,s_{2,r}}}{\tau_r}
                   +\frac{n\ell_r}{\tau_r}\right]\tau_r
 \notag\\
 &\le C_{L,\delta}
       [1+\sqrt{A\rho^{-1}H}+AH]\tau_r,
       \label{eq:sparse-coarse-gram}\\
 \|(Q_u)_{SS}-\tau_uI_n\|_{\mathrm{op}}
 &\le\bigl[C_{L,\delta}
       (\sqrt{A\rho H}+A\rho^2H)+o(1)\bigr]\tau_r.
       \label{eq:sparse-fine-gram}
\end{align}
The $o(1)$ is uniform over the stated range. In particular, after
reducing $\alpha_0$ if necessary,
$\|(Q_u)_{SS}\|_{\mathrm{op}}\le2\tau_r$ on the same event for all
sufficiently large $d$. Deterministically,
\begin{align}
 0\le\tau_r-\tau_u&\le C_L\tau_r/r,
       \label{eq:gram-sparse-band-trace}\\
 n\ell_u/\tau_r&\le C_LA\rho^2H,
       \label{eq:sparse-fine-eigenvalue}
\end{align}
where $\ell_v=\lambda_{v+1}$. No simultaneous event over different
sample sizes is asserted.
\end{lemma}
\begin{proof}
The dimension and spectral estimates already proved give
\begin{align*}
 n&\le C A\binom D{r+2}rH,&
 \log(2n)&\le C_A rH,\\
 n\ell_r/\tau_r&\le C_LAH,&
 nd\,s_{2,r}/\tau_r^2&\le C_LA\rho^{-1}H.
\end{align*}
For $u=r+2$, the exact two-step binomial ratio gives
\begin{align*}
 n\ell_u/\tau_r&\le C_LA\rho^2H,&
 nd\,s_{2,u}/\tau_r^2&\le C_LA\rho H.
\end{align*}
Moreover $\tau_r-\tau_u=\binom D{r+2}a_{r+2}$, so the level and
trace estimates prove \eqref{eq:gram-sparse-band-trace}.

We apply the proof of Lemma~\ref{lem:gram-general} with a different
deterministic mask. Retain pairs satisfying
\[
 \left|\frac{1+x\cdot y}{d+1}\right|\le a_n,
 \qquad a_n=\sqrt{128\log(2n)/d}.
\]
By choosing $\alpha_0$ sufficiently small, uniformly
$\log(2n)/d\le C_A\alpha\log(e/\alpha)+o(1)$ is as small as needed,
so $a_n\le a_0$. For independent uniform inputs, Hoeffding's bound
(and $a_n\ge2/(d+1)$) gives exceptional-pair probability
\[
 b_n\le2\exp(-a_n^2d/8)=2(2n)^{-16}.
\]
The exceptional set is translation invariant on the cube. Its size
is therefore the same in every row, so the masked-matrix proof applies
with $B_{\mathrm{near}}/M=b_n$. This mask is deterministic for each
fixed pair $(d,n)$; it does not depend on the observed sample.
Lemma~\ref{lem:near-orthogonal} supplies the retained-entry bound
\begin{equation}\label{eq:sparse-coherence}
 \frac{\eta_v}{\tau_v}
 \le\left[C\sqrt{\frac{\log(2n)+v+2}{d}}\right]^{v+2}
       +C d^\beta e^{-cd},\qquad v\in\{r,u\}.
\end{equation}
Thus exactly the argument of Lemma~\ref{lem:gram-general}, still with
$q=2\log M=2d\log2$, gives for every $z>0$
\begin{align}
 &\mathbb P\{\|(Q_v)_{SS}-\tau_vI_n\|>z\tau_v\}\notag\\
 &\quad\le\frac C{z^2}
 \left[\frac{q\eta_v}{\tau_v}
       +\frac{\sqrt{nq s_{2,v}}}{\tau_v}
       +\frac{n\ell_v}{\tau_v}+n(M^{-1}+b_n)\right]^2
       +n^2(M^{-1}+b_n).
 \label{eq:sparse-masked-tropp}
\end{align}
The collision and exceptional-pair terms tend to zero uniformly:
$n\ge A D_{r_0}\log(D_{r_0}+1)\to\infty$, and $\alpha_0$ can be chosen
so that $n\le e^{c_1d}$ with $2c_1<\log2$.

The coherence term $q\eta_v/\tau_v$ also tends to zero uniformly.
For $r\le\sqrt d$, its bracket in \eqref{eq:sparse-coherence} is at
most $C_A d^{-1/4}\sqrt{\log d}$, and its exponent is at least
$r_0+2\ge13$. Hence it is bounded, after multiplication by $q$, by
\[
 C d[C_A d^{-1/4}\sqrt{\log d}]^{13}
       +C d^{1+\beta}e^{-cd}=o(1).
\]
For $r>\sqrt d$, reducing $\alpha_0$ makes the bracket at most $1/2$,
so the bound is $Cd\,2^{-\sqrt d}+Cd^{1+\beta}e^{-cd}=o(1)$.
Using \eqref{eq:sparse-masked-tropp} for $v=r,u$, with failure
probability at most $\delta/2$ for each, now proves
\eqref{eq:sparse-coarse-gram} and \eqref{eq:sparse-fine-gram}.
For the latter, multiply by $\tau_u/\tau_r\le1$ before using the
spectral bounds, so no lower comparison of the two traces is needed.
Finally $\rho H\to0$ as $\rho\to0$, uniformly for $0<\rho\le\alpha$;
choose $\alpha_0$ to make the fixed part of the fine deviation below
$1/2$, and absorb the uniform $o(1)$. Since $\tau_u\le\tau_r$, this
gives the asserted upper bound $2\tau_r$.
\end{proof}

% Standalone section: exact positive-subcritical posterior lower bound.
\section{An exponential sample lower bound}
\label{sec:positive-sparse-partial}
Throughout this section $f=\chi_U$, where $U\subseteq[d]$ and $|U|=k$.
Write $D=d+1$, $\beta=2^{-(L-1)}$, and
\[
 B_{d,k}=\sqrt{k}\log\frac{eD}{\sqrt{k}},\qquad
 W_{d,k}=\max\{k,B_{d,k}\}.
\]
The $k$ term follows from the random-restriction amplitude tradeoff.
The second term is useful in very large ambient dimensions and follows
from a direct polynomial-threshold obstruction.

\begin{theorem}[Sparse parity for $0<\gamma<1$]
\label{thm:positive-sparse-partial}
Fix $L\ge2$ and $a>0$. There are constants $c>0$, $k_0$, and $d_0$,
depending only on $L,a$, such that the following holds. If
$d\ge d_0$, $k_0\le k\le d$, and a specified sample size satisfies
\[
 0\le n\le\exp(cW_{d,k}),
\]
then, with probability at least $3/4$ over independent uniform training
inputs, every minimizer of
\[
 \min\left\{\frac12\|z\|_K^2:
 |z(x)|\ge a\ \text{for all }x,\quad
 f(X_i)z(X_i)\ge a\ \text{for all }i\right\}
\]
has population sign error at least $3/8$.
On the same event, for every fixed $0<\gamma<1$, the actual network
posterior in the width-first limit satisfies
\[
 \liminf_{N\to\infty}
 \mathbb P_x\{q_{N,S}(x)\le2/3\}\ge\frac1{16}.
\]
The probability statement is for each specified sample size; it is not
a simultaneous event over all sample prefixes.
\end{theorem}

In particular, the lower bound is exponential in $k$ throughout the
entire growing sparse regime. It does not give the stronger scale
$\exp\{c k\log(ed/k)\}$ when $d/k\to\infty$.

\subsection{Squared balance on the target coordinates}
\begin{lemma}[Active-coordinate squared balance]
\label{lem:positive-active}
If $P$ has total Walsh degree at most $r$, then
\begin{equation}
 \mathbb E(fP)_-^2
 \ge\left(\frac12-\sqrt{\frac rk}\right)\|P\|_2^2.
 \label{eq:positive-active-balance}
\end{equation}
For every fixed $\delta>0$, there is a universal $A_\delta$ such that
if $r\le k/64$ and $n\ge A_\delta D_r\log(D_r+1)$, then, with
probability at least $1-\delta$,
\begin{equation}
 \mathbb P_n(fP)_-^2\ge\frac14\|P\|_2^2
 \quad\text{for every polynomial $P$ of degree at most $r$}.
 \label{eq:positive-active-sampling}
\end{equation}
\end{lemma}
\begin{proof}
For $j\in U$ put $D_jP=(P-P\circ\mathrm{flip}_j)/2$.
Because $f$ changes sign under these flips, the inequality
$|F(s)-F(t)|\le(|s|+|t|)|s-t|$ for $F(t)=t|t|$ gives
\[
 |\mathbb E fP|P||\le2\|P\|_2\|D_jP\|_2.
\]
Averaging over $j\in U$ and using
$\sum_{j\in U}\|D_jP\|_2^2\le r\|P\|_2^2$ proves
\eqref{eq:positive-active-balance} from
$(fP)_-^2=(P^2-fP|P|)/2$.
The vector of all Walsh characters of degree at most $r$ is isotropic
and has squared length $D_r$ at every input. Under $r\le k/64$,
the population bound is at least $3\|P\|_2^2/8$.
Apply Lemma~\ref{lem:squared-sampling} at accuracy $1/8$ and use
homogeneity to obtain \eqref{eq:positive-active-sampling}.
\end{proof}

\subsection{The polynomial-threshold range}
\begin{lemma}[A $\sqrt{k}$ cutoff]
\label{lem:positive-sqrt-range}
For fixed $L,a$, there are $c_b>0$, $k_b$, and $d_b$ such that, on the
localized-comparator event, every minimizer has error at least $3/8$
whenever $d\ge d_b$, $k_b\le k\le d$, and
\[
 n\le\exp\left(c_b\sqrt{k}\log\frac{eD}{\sqrt{k}}\right).
\]
\end{lemma}
\begin{proof}
The comparator gives the simultaneous bound
\begin{equation}
 \|z\|_K^2\le C_{L,a}(d+n d^\beta)
 \label{eq:positive-comparator-budget}
\end{equation}
for every minimizer. Choose a sufficiently small constant $\kappa>0$
and an odd $s$ with $s\sim\kappa\sqrt{k}$.
Fixing the $d-k$ nuisance coordinates leaves a degree-$s$ polynomial
on the active $k$-cube. Lemma~\ref{lem:ptf-parity} and the central
binomial bound imply that every such polynomial threshold has error
at least $7/16$, for all sufficiently large $k$.

The sharp spectrum gives
\[
 \|z-z_{\le s}\|_2^2
 \le C_{L,a}(d+n d^\beta)\lambda_{s+1}
 \le\frac{C_{L,a}(d+n d^\beta)s^{-1-\beta}}
            {\binom D{s+2}}.
\]
The binomial lower bound shows
\[
 \log\binom D{s+2}
 \ge c_1\kappa\sqrt{k}\log\frac{eD}{\sqrt{k}}
\]
for a universal $c_1>0$ once $k$ is large.
Choose $c_b<c_1\kappa/4$, then increase $k_b$ to absorb the factor
$d^\beta$ and increase $d_b$ to make the displayed tail bound at most
$a^2/16$, uniformly in the stated range. The additive $d$ is harmless
because $s$ is then a sufficiently large fixed integer or larger.
Since $|z|\ge a$, its sign differs from a fixed-tie threshold of
$z_{\le s}$ on at most $1/16$ of the cube. Its error is therefore
at least $7/16-1/16=3/8$.
The comparator's allowed sample range contains the present one for all
large $d$, because $\sqrt{k}\log(eD/\sqrt{k})=o(d)$ uniformly in $k\le d$.
\end{proof}

\subsection{Proof of the positive-subcritical theorem}
\begin{proof}[Proof of Theorem~\ref{thm:positive-sparse-partial}]
Fix $\delta=1/40$, choose $A$ large enough for
Lemma~\ref{lem:positive-active}, and fix an odd $r_0\ge11$.
Choose the small degree fraction in
Lemma~\ref{lem:sparse-gram} after these choices.
The localized comparator has failure probability $e^{-\Omega_L(d)}$.
All its conclusions hold simultaneously over the minimizing fields.

If $W_{d,k}=B_{d,k}$, choose $c\le c_b$ and apply
Lemma~\ref{lem:positive-sqrt-range}. We may therefore assume
$W_{d,k}=k$ and $n\le e^{ck}$.
Let $r$ be the largest odd integer satisfying
\[
 A D_r\log(D_r+1)\le n.
\]
If this cutoff does not exist or $r<r_0$, then
$n<A D_{r_0}\log(D_{r_0}+1)$.
After increasing $k_0$, this entire range lies below the threshold of
Lemma~\ref{lem:positive-sqrt-range}, which again proves the claim.

Suppose $r\ge r_0$. Choose $c<(\log2)/4$.
Since $k\le d$ and $D_r\ge2^{d-1}$ for $r\ge d/2$, the defining
lower bound on $n$ excludes $r\ge d/2$ for sufficiently large $d$.
In this range, the binomial lower bound implies
\begin{equation}
 rH\le C_0\log n\le C_0 c k,
 \qquad \rho=\frac rD,\quad H=\log\frac e\rho,
 \label{eq:positive-main-cutoff}
\end{equation}
for a universal constant $C_0$.
By taking $c$ sufficiently small, we ensure $r\le k/64$,
$r\le\alpha d$, and $n\le e^{d/25}$.
All comparator and sparse-Gram hypotheses therefore hold.
Moreover $r\ge r_0$ makes $n$ polynomially larger than $d$, so
\eqref{eq:positive-comparator-budget} becomes
\begin{equation}
 B:=\|z\|_K^2\le C_{L,a}n d^\beta.
 \label{eq:positive-main-budget}
\end{equation}

Work on the intersection of the comparator event, the squared-sampling
event \eqref{eq:positive-active-sampling}, and the Gram event in
Lemma~\ref{lem:sparse-gram}. Their total failure probability is
at most $2\delta+e^{-\Omega_L(d)}<1/4$ for sufficiently large $d$.
For an arbitrary minimizer write
\[
 z=P+b+t,\qquad P=z_{\le r},\quad b=z_{(r,r+2]},\quad t=z_{>r+2}.
\]
At a training point where $fP<0$, consistency gives
$(fP)_-\le|z-P|$. The sampling bound, RKHS evaluation bound, and
\eqref{eq:sparse-coarse-gram} yield
\[
 \frac14\|P\|_2^2
 \le\mathbb P_n(fP)_-^2
 \le\mathbb P_n(z-P)^2
 \le\frac Bn\|(Q_r)_{SS}\|_{\mathrm{op}}.
\]
Using \eqref{eq:positive-main-budget}, $\tau_r\asymp r^{-\beta}$,
and the deterministic estimates
$n\lambda_{r+1}\le C_L A H\tau_r$ and
$n\lambda_{r+3}\le C_L A\rho^2H\tau_r$, we obtain, simultaneously
for every minimizer,
\begin{align}
 \|P\|_2^2&\le C_{L,a,A,\delta}\rho^{-(\beta+1)}H,
 \label{eq:positive-partial-head}\\
 \|b\|_2^2&\le B\lambda_{r+1}
             \le C_{L,a,A}\rho^{-\beta}H,
 \label{eq:positive-partial-band}\\
 \|t\|_2^2&\le B\lambda_{r+3}
             \le C_{L,a,A}\rho^{2-\beta}H.
 \label{eq:positive-partial-tail}
\end{align}
For the first estimate we bounded
$1+\sqrt{A\rho^{-1}H}+AH$ by $C_A\rho^{-1}H$.
The eigenvalue estimates follow from the sharp spectrum and the
maximality interval for $r$; the second is also stated in
\eqref{eq:sparse-fine-eigenvalue}.
Put $h=P+b=z_{\le u}$, where $u=r+2\le2r$.
The first two pieces are orthogonal in $L^2$, and their sum therefore
has squared norm at most $C\rho^{-(\beta+1)}H$. Hence
Lemma~\ref{cor:restriction-amplitude} gives
\begin{equation}
 \mathbb P\{\operatorname{sign}z\ne f\}
 \ge\frac7{16}
       -C_1\rho^{2-\beta}H
       -C_1\rho^{-(\beta+1)}H\exp\!\left(-c_*\frac{k}{u}\right),
 \label{eq:positive-restriction-risk}
\end{equation}
where $c_*>0$ is universal and $C_1$ depends only on the already fixed
constants $L,a,A,\delta$.

We now choose the theorem's $c$. From
\eqref{eq:positive-main-cutoff}, $\rho\le C_0c$ and
$k/u\ge H/(2C_0c)$.
Since $x^{2-\beta}\log(e/x)\to0$ as $x\downarrow0$, the first error
term in \eqref{eq:positive-restriction-risk} is at most $1/32$ once
$c$ is sufficiently small. For the second term, use
$\rho^{-(\beta+1)}=\exp[(\beta+1)(H-1)]$ to obtain
\[
 C_1\rho^{-(\beta+1)}H e^{-c_*k/u}
 \le C_1 H\exp\!\left[
       \left(\beta+1-\frac{c_*}{2C_0c}\right)H\right].
\]
As $H\ge1$, taking $c$ still smaller makes this at most $1/32$
uniformly. These choices also ensure all earlier degree-fraction
restrictions. Finally increase $k_0,d_0$ to absorb finite rounding and
uniform asymptotic estimates. Thus every minimizing sign has error at
least $7/16-1/32-1/32=3/8$.

For completeness, let $q(x)$ be any probability mixture of these
minimizing signs, let $R=\mathbb E_x(1-q(x))$, and put
$F=\mathbb P_x\{q(x)\le2/3\}$.
The draw-error bound gives $R\ge3/8$, while
\[
 R\le\frac{1-F}{3}+F,
 \qquad F\ge\frac{3R-1}{2}\ge\frac1{16}.
\]
The positive-gamma width-limit proposition says that, for each fixed
$0<\gamma<1$, the actual finite-width posterior puts mass $1-o(1)$
on these minimum-energy signs. Thus its draw error is at least
$3/8-o(1)$, independently of the relative weights among ties.
Apply the same inequality at every width and take the lower limit.
This proves the asserted actual-posterior guarantee.
\end{proof}

\section{An upper bound and the dependence on ambient dimension}
\label{sec:upper-bound}
An upper bound identifies the natural comparison scale in the sparse
regime.

\begin{proposition}[An actual-posterior upper bound]\label{prop:upper-bound}
Let $\lambda_k$ be the normalized Walsh eigenvalue defined above.
There is an absolute constant $C$ such that, for every
$0<\varepsilon<1/3$, $0<\delta<1$, and
\[
 n\ge \frac{C}{\varepsilon^2}
       \left(\lambda_k^{-1}+\log\frac2\delta\right),
\]
the actual posterior at every fixed $0<\gamma<1$ satisfies
\[
 \mathbb P_S\left\{\limsup_{N\to\infty}F_N(S)
                      \le3\varepsilon\right\}\ge1-\delta.
\]
On the same sample event, every compatible minimum-energy sign function
has test error at most $\varepsilon$.
\end{proposition}
\begin{proof}
The globally separated comparator $z=af$ is feasible for every sample
and has $\|af\|_K^2=a^2/\lambda_k$. Thus every minimizing field lies
in the RKHS ball of radius $R=a/\sqrt{\lambda_k}$ and satisfies
$f(X_i)z(X_i)\ge a$.
Use the ramp loss
$\ell(t)=\min\{1,\max\{0,1-t/a\}\}$.
Its empirical loss vanishes for these fields, and
$\mathbf1_{\{f(x)z(x)\le0\}}\le\ell(f(x)z(x))$.

For completeness, the standard uniform deviation estimate here is
\[
 \sup_{\|z\|_K\le R}
 \left[\mathbb E\ell(fz)-\frac1n\sum_i\ell(f(X_i)z(X_i))\right]
 \le \frac{4R}{a\sqrt n}
       +\sqrt{\frac{\log(1/\delta)}{2n}}
\]
with probability at least $1-\delta$.
Indeed, symmetrization and contraction for the $a^{-1}$-Lipschitz
function $\ell-\ell(0)$ bound the expected supremum by $4/a$ times
the RKHS Rademacher complexity. The reproducing identity and
$K(x,x)=1$ give
\[
 \frac1n\mathbb E_\sigma
 \sup_{\|z\|_K\le R}\left|\sum_i\sigma_i f(X_i)z(X_i)\right|
 =\frac Rn\mathbb E_\sigma
       \left\|\sum_i\sigma_i f(X_i)K_{X_i}\right\|_K
 \le\frac R{\sqrt n}.
\]
Changing one sample changes the supremum deviation by at most $1/n$;
the bounded-differences inequality supplies the displayed probability
bound. The assumed sample size therefore makes every minimizing
sign's error at most $\varepsilon$.

By Proposition~\ref{prop:rate}, the actual posterior mass outside these
minimizing signs tends to zero. Hence its average test error
$\mathbb E_x[1-q_{N,S}(x)]$ has upper limit at most $\varepsilon$.
Since $F_N(S)\le3\mathbb E_x[1-q_{N,S}(x)]$, the claimed confidence
bound follows, including the strict greater-than-$2/3$ convention.
\end{proof}

For growing $k=o(d)$, put $j=k$ if $k$ is odd and $j=k+1$ otherwise.
Lemma~\ref{lem:sharp} gives
\[
 \lambda_k^{-1}\asymp j^{1+\beta}\binom{d+1}{j},
 \qquad
 \log\lambda_k^{-1}=\Theta\!\left(k\log\frac{ed}{k}\right).
\]
In particular, when $k=\lfloor d/\log d\rfloor$, fixed
$\varepsilon,\delta>0$ require only
$\exp\{O(k\log\log d)\}=d^{o(k)}$ samples in this upper bound.
Thus a literal $d^{\Omega(k)}$ lower bound, with a fixed positive
exponent uniformly over all growing sparse $k=o(d)$, is false for the
actual posterior. The upper bound has the natural exponential scale
$\exp\{\Theta(k\log(ed/k))\}$. A lower bound on this scale remains
open. For $k\le d^{1-\eta}$ with fixed $\eta>0$, such a lower bound
would entail $d^{\Omega(k)}$.

\section{Why an amplitude bound alone cannot settle the sparse case}
\label{sec:blocked-obstruction}
The remaining difficulty cannot be removed simply by proving that every
accurate low-degree threshold approximation to parity has amplitude growing
with the number of active coordinates. The following construction rules
out that assertion without any symmetry assumption on the polynomial.

\begin{proposition}[Bounded-amplitude blocked approximants]
\label{prop:blocked-approximants}
For every $\alpha,\varepsilon>0$, there are finite constants
$B=B(\alpha,\varepsilon)$ and $k_0=k_0(\alpha,\varepsilon)$ such that,
for every $d\ge k\ge k_0$, a multilinear polynomial $P$ on the $d$-cube
satisfies
\[
 \deg P\le\alpha k,\qquad
 1\le |P(x)|\le B\quad\hbox{for every }x,\qquad
 \Pr\{\operatorname{sign}P(X)\ne\chi_U(X)\}\le\varepsilon
\]
for any fixed $U\subseteq[d]$ of size $k$.
\end{proposition}
\begin{proof}
Replacing $\alpha,\varepsilon$ by smaller positive constants if necessary,
assume $\alpha\le1$ and $\varepsilon\le1/2$.
Choose a fixed integer $m$ sufficiently large that a central integer
interval $[A,B_0]$ for $T\sim\operatorname{Bin}(m,1/2)$ has probability
at least $1-\varepsilon$ and length $s=B_0-A\le\alpha m/2$.
Hoeffding's inequality permits $m=O(\alpha^{-2}\log(2/\varepsilon))$,
after increasing the absolute constant. Define
\[
 Q(t)=\sigma\prod_{j=A}^{B_0-1}(2t-2j-1),
\]
where $\sigma\in\{-1,1\}$ makes its sign agree with $(-1)^t$ throughout
$[A,B_0]\cap\mathbb Z$. At every integer $t\in[0,m]$, all factors are
nonzero odd integers, so $1\le|Q(t)|\le(2m+1)^s$.
Partition $U$ into $m$ nonempty blocks of size at most $\lceil k/m\rceil$
and let $Y_j$ be the parity of the $j$th block. These are independent
uniform signs, and their product is $\chi_U$. Set
\[
 P(x)=Q\!\left(\frac{m-\sum_{j=1}^mY_j(x)}2\right).
\]
Multilinearization in the block variables shows
$\deg P\le s\lceil k/m\rceil\le\alpha k$ whenever $k\ge m$.
The remaining conclusions follow from the bounds on $Q$ and the
binomial law of its argument.
\end{proof}

For example, if $3$ divides $k$, three equal blocks give
$P=-(Y_1+Y_2+Y_3)$, with degree $k/3$, global margin one, parity error
$1/4$, and $\|P\|_2^2=3$. If $4$ divides $k$, the polynomial
$P=1-\sum_{j=1}^4Y_j$ has degree $k/4$, error $5/16$, and
$\|P\|_2^2=5$. These examples need not be symmetric in all $k$ original
coordinates.

Proposition~\ref{prop:blocked-approximants} is an obstruction to a proof
method, not a learning upper bound for the posterior. If $\deg P<k$,
then $\mathbb E[\chi_U P]=0$; together with a global unit margin and
sign error $R$, this implies
\[
 \mathbb E(1-\chi_U P)_+^2\ge\frac1R.
\]
Thus an accurate blocked approximant has larger squared repair loss
than the constant polynomial, whose loss is exactly two. Determining
the minimum energy of repairs that also preserve the global margin
remains essential.

\section{A small bounded-weight parity network}\label{sec:paritycircuit}
For completeness, full parity on $k$ bits has a one-hidden-layer
$\tanh$ realization of width $k+1$, with weights $O(k\log(k+1))$
and output margin greater than $1/2$. Let
\[
 t(x)=\frac{k+\sum_i x_i}{2}\in\{0,\ldots,k\},\quad
 y_j=(-1)^{k-j},\quad A=\log(16(k+1)).
\]
The step functions $s_j(t)=\sign(t-j+1/2)$ satisfy
\[
 y_t=c+\sum_{j=1}^k y_js_j(t),\qquad c=(y_0+y_k)/2\in\{0,1\}.
\]
Replace each step by $h_j=\tanh(A(t-j+1/2))$. Since its argument
has absolute value at least $A/2$, its error is at most $2e^{-A}$.
A constant neuron $h_0=\tanh1$ supplies the term $c$ using output
weight $c/\tanh1$. The resulting output differs from parity by
less than $1/8$. First-layer input weights are $A/2$ and bias
weights have magnitude at most $Ak/2$; readout weights have
magnitude at most $2$.

At depth $L\ge3$, the next hidden layer sums these features and
applies $\tanh$, producing the correct sign with magnitude at least
$\tanh(7/8)>1/2$. Each additional hidden layer can preserve this
margin by a single weight $2$, since $\tanh(2u)>1/2$ whenever
$u>1/2$. A final readout weight $1$ completes the construction.
Unused neurons are set to zero. Thus the same polynomial bounds
hold at every fixed depth.
For any other fixed margin $a>0$, multiplying the final readout by $2a$
gives margin greater than $a$ with the same polynomial complexity bounds.



\section{Scope}
Theorem~\ref{thm:positive-sparse-partial}, together with the actual-network
rate in Proposition~\ref{prop:rate}, proves Theorem~\ref{thm:main}.
Parity has a Boolean circuit with $O(k)$ gates and the bounded-weight
neural realization given above. Its difficulty here is a property of the
specified posterior. An algorithm that searches parity supports can
identify the target with $O(k\log(ed/k))$ noiseless examples at fixed
success probability: each incorrect parity agrees with the target on
each independent uniform example with probability $1/2$, and a union
bound applies to the candidate supports.

Section~\ref{sec:upper-bound} gives the corresponding actual-posterior
upper bound and explains the dependence on ambient dimension. A matching
$\exp\{c k\log(ed/k)\}$ lower bound remains open in the growing sparse
regime, including a $d^{\Omega(k)}$ lower bound when $k\le d^{1-\eta}$
for fixed $\eta>0$. No assertion about the endpoint $\gamma=0$ is made
in this paper.

\begin{thebibliography}{9}
\bibitem{ABFR}
J. Aspnes, R. Beigel, M. Furst, and S. Rudich.
The expressive power of voting polynomials.
\emph{Combinatorica} 14(2):135--148, 1994.
\href{https://www.cs.yale.edu/homes/aspnes/papers/stoc91voting.pdf}{Author manuscript}.
\bibitem{flajolet2009}
P. Flajolet and R. Sedgewick. \emph{Analytic Combinatorics}.
Cambridge University Press, 2009. Chapter VI.
\bibitem{Lee2023}
J. R. Lee. Spectral hypergraph sparsification via chaining.
\href{https://arxiv.org/abs/2209.04539}{arXiv:2209.04539},
Section 2.4, Lemma 2.12 and equation (2.20).
\bibitem{tropp2008}
J. A. Tropp. Norms of random submatrices and sparse approximation.
\emph{Comptes Rendus Math\'ematique} 346:1271--1274, 2008.
\href{https://tropp.caltech.edu/papers/Tro08-Norms-Random-v3.pdf}{Author manuscript},
Theorem 1.1.
\end{thebibliography}
\end{document}
